% GENERATED FILE. Edit canonical lessons, then run npm run generate:books.
% canonical-source-hash: fnv1a64:53f0628eb915329f
% canonical-lessons: AR-C138-nasr, AR-C138-difda, AR-C138-khadd, AR-C138-jabin, AR-C138-hajib

\chapter{Eagle, Frog, Cheek, Forehead, Eyebrow}
\label{ch:ar-eagle-frog-cheek-forehead-eyebrow}
\hlchaptermodality{138}

\begin{chapteropening}
\textbf{By the end of this chapter:} \emph{I can say an eagle, a frog, a cheek, a forehead and an eyebrow.}

Complete the last lesson of chapter 138: I can say an eagle, a frog, a cheek, a forehead and an eyebrow.
\end{chapteropening}

\section[\texorpdfstring{\ar{نسر} (nasr)}{نسر (nasr)}]{\ar{نسر} (nasr) — an eagle}

\label{lesson:AR-C138-nasr}



\begin{quote}
Before the new one: say the Arabic for an owl, then the Arabic for a crow.
\end{quote}

\subsection*{You'll want to know: \ar{نسر}}
\textbf{\ar{نسر}} — \emph{nasr} — \textquotedblleft{}an eagle\textquotedblright{}.

\begin{cousinweb}
\textbf{\ar{نسر}} (\emph{nasr}) is built on the native Arabic root \emph{n-s-r}, \textquotedblleft{}a vulture, an eagle\textquotedblright{} (Lane). Lane compares Hebrew \emph{nesher}, \textquotedblleft{}eagle\textquotedblright{}.
\end{cousinweb}

\begin{grammarlens}[title={What you've built}]
One more word for this chapter.
\end{grammarlens}

\subsection*{Your turn}
\begin{itemize}
\raggedright
  \item \emph{Say it:} \emph{nasr}
  \item \emph{Say it:} \emph{nasr}, once more
  \item \emph{Say it:} say \emph{nasr}
  \item \emph{Recall:} say the Arabic for an owl, then the Arabic for a crow, then say \emph{nasr} again
\end{itemize}

\begin{tcolorbox}[breakable,colback=teal!4,colframe=teal!35!black,title={Before you move on}]
What does \ar{نسر} mean? (\textquotedblleft{}An eagle\textquotedblright{}.) Say it once more.
\end{tcolorbox}
\section[\texorpdfstring{\ar{ضفدع} (ḍifdaʿ)}{ضفدع (ḍifdaʿ)}]{\ar{ضفدع} (ḍifdaʿ) — a frog}

\label{lesson:AR-C138-difda}



\begin{quote}
Before the new one: say the Arabic for a crow, then the Arabic for an eagle.
\end{quote}

\subsection*{You'll want to know: \ar{ضفدع}}
\textbf{\ar{ضفدع}} — \emph{ḍifdaʿ} — \textquotedblleft{}a frog\textquotedblright{}.

\begin{cousinweb}
\textbf{\ar{ضفدع}} (\emph{ḍifdaʿ}) is built on the native Arabic root \emph{ḍ-f-d-ʿ}, \textquotedblleft{}a frog\textquotedblright{} (Lane). Hebrew \emph{tsefardeaʿ}, \textquotedblleft{}frog\textquotedblright{}, is usually counted a cousin.
\end{cousinweb}

\begin{grammarlens}[title={What you've built}]
One more word for this chapter.
\end{grammarlens}

\subsection*{Your turn}
\begin{itemize}
\raggedright
  \item \emph{Say it:} \emph{ḍifdaʿ}
  \item \emph{Say it:} \emph{ḍifdaʿ}, once more
  \item \emph{Say it:} say \emph{ḍifdaʿ}
  \item \emph{Recall:} say the Arabic for a crow, then the Arabic for an eagle, then say \emph{ḍifdaʿ} again
\end{itemize}

\begin{tcolorbox}[breakable,colback=teal!4,colframe=teal!35!black,title={Before you move on}]
What does \ar{ضفدع} mean? (\textquotedblleft{}A frog\textquotedblright{}.) Say it once more.
\end{tcolorbox}
\section[\texorpdfstring{\ar{خد} (khadd)}{خد (khadd)}]{\ar{خد} (khadd) — a cheek}

\label{lesson:AR-C138-khadd}



\begin{quote}
Before the new one: say the Arabic for an eagle, then the Arabic for a frog.

Greet the morning, and give its answer. (\emph{Ṣabāḥ al-khayr}; \emph{ṣabāḥ an-nūr}.)
\end{quote}

\subsection*{You'll want to know: \ar{خد}}
\textbf{\ar{خد}} — \emph{khadd} — \textquotedblleft{}a cheek\textquotedblright{}.

\begin{cousinweb}
\textbf{\ar{خد}} (\emph{khadd}) is built on the native Arabic root \emph{kh-d-d}, \textquotedblleft{}the cheek\textquotedblright{} (Lane). An old noun with its own root in Lane.
\end{cousinweb}

\begin{grammarlens}[title={What you've built}]
One more word for this chapter.
\end{grammarlens}

\subsection*{Your turn}
\begin{itemize}
\raggedright
  \item \emph{Say it:} \emph{khadd}
  \item \emph{Say it:} \emph{khadd}, once more
  \item \emph{Say it:} say \emph{khadd}
  \item \emph{Recall:} say the Arabic for an eagle, then the Arabic for a frog, then say \emph{khadd} again
\end{itemize}

\begin{tcolorbox}[breakable,colback=teal!4,colframe=teal!35!black,title={Before you move on}]
What does \ar{خد} mean? (\textquotedblleft{}A cheek\textquotedblright{}.) Say it once more.
\end{tcolorbox}
\section[\texorpdfstring{\ar{جبين} (jabīn)}{جبين (jabīn)}]{\ar{جبين} (jabīn) — a forehead}

\label{lesson:AR-C138-jabin}



\begin{quote}
Before the new one: say the Arabic for a frog, then the Arabic for a cheek.
\end{quote}

\subsection*{You'll want to know: \ar{جبين}}
\textbf{\ar{جبين}} — \emph{jabīn} — \textquotedblleft{}a forehead\textquotedblright{}.

\begin{cousinweb}
\textbf{\ar{جبين}} (\emph{jabīn}) is built on the native Arabic root \emph{j-b-n}, \textquotedblleft{}the side of the forehead\textquotedblright{} (Lane). Lane's \emph{jabīn} is the side of the forehead above the temple.
\end{cousinweb}

\begin{grammarlens}[title={What you've built}]
One more word for this chapter.
\end{grammarlens}

\subsection*{Your turn}
\begin{itemize}
\raggedright
  \item \emph{Say it:} \emph{jabīn}
  \item \emph{Say it:} \emph{jabīn}, once more
  \item \emph{Say it:} say \emph{jabīn}
  \item \emph{Recall:} say the Arabic for a frog, then the Arabic for a cheek, then say \emph{jabīn} again
\end{itemize}

\begin{tcolorbox}[breakable,colback=teal!4,colframe=teal!35!black,title={Before you move on}]
What does \ar{جبين} mean? (\textquotedblleft{}A forehead\textquotedblright{}.) Say it once more.
\end{tcolorbox}
\section[\texorpdfstring{\ar{حاجب} (ḥājib)}{حاجب (ḥājib)}]{\ar{حاجب} (ḥājib) — an eyebrow}

\label{lesson:AR-C138-hajib}



\begin{quote}
Before the new one: say the Arabic for a cheek, then the Arabic for a forehead.
\end{quote}

\subsection*{You'll want to know: \ar{حاجب}}
\textbf{\ar{حاجب}} — \emph{ḥājib} — \textquotedblleft{}an eyebrow\textquotedblright{}.

\begin{cousinweb}
\textbf{\ar{حاجب}} (\emph{ḥājib}) is built on the native Arabic root \emph{ḥ-j-b}, \textquotedblleft{}to veil, to screen\textquotedblright{} (Lane). \emph{Ḥājib} is \textquotedblleft{}that which screens\textquotedblright{}: the eyebrow, and a chamberlain; the same root gives \emph{ḥijāb}.
\end{cousinweb}

\begin{grammarlens}[title={What you've built}]
One more word for this chapter.
\end{grammarlens}

\subsection*{Your turn}
\begin{itemize}
\raggedright
  \item \emph{Say it:} \emph{ḥājib}
  \item \emph{Say it:} \emph{ḥājib}, once more
  \item \emph{Say it:} say \emph{ḥājib}
  \item \emph{Recall:} say the Arabic for a cheek, then the Arabic for a forehead, then say \emph{ḥājib} again
\end{itemize}

\begin{tcolorbox}[breakable,colback=teal!4,colframe=teal!35!black,title={Before you move on}]
What does \ar{حاجب} mean? (\textquotedblleft{}An eyebrow\textquotedblright{}.) Say it once more.
\end{tcolorbox}
